% Compile with: uplatex file.tex && dvipdfmx file.dvi
%           or: latexmk -uplatex file.tex
\documentclass[dvipdfmx]{article}
\usepackage{amsmath, amssymb}
\usepackage{tikz}

\title{scenario\_fourier\_convolution}
\date{\today}

\begin{document}
\maketitle

Write the Fourier transform of the Gaussian $e^{-t^2}$ as the definition's own definite integral, $F(\omega)=\int_{-\infty}^{\infty} f(t)e^{-i\omega t}dt$. No dedicated operation is used.

$F = \int_{-\infty}^{\infty} e^{-t^2} e^{-i \omega t} dt$

\begin{equation}
  F = \int_{-\infty}^{\infty} e^{-t^2} e^{-i \omega t} dt
  \label{eq:E1}
\end{equation}

Evaluate the right-hand side.

compute the right-hand side of E1

\begin{equation}
  F = e^ {- \left(\frac{\omega^2}{4}\right) } \sqrt{\pi}
  \label{eq:E2}
\end{equation}

Write the inverse from its definition too, $\frac{1}{2\pi}\int F(\omega)e^{i\omega t}d\omega$.

$f = \frac{1}{2\pi}\int_{-\infty}^{\infty} \sqrt{\pi} e^{-\frac{\omega^2}{4}} e^{i \omega t} d\omega$

\begin{equation}
  f = \frac{1}{2\pi}\int_{-\infty}^{\infty} \sqrt{\pi} e^{-\frac{\omega^2}{4}} e^{i \omega t} d\omega
  \label{eq:E3}
\end{equation}

Evaluating returns exactly the original $e^{-t^2}$ — the round trip closes.

compute the right-hand side of E3

\begin{equation}
  f = e^ {- t^2 }
  \label{eq:E4}
\end{equation}

Before the parametrised Gaussian, state the convergence condition $a>0$ explicitly. The macro silently assumed every symbolic parameter in the expression positive, but that is a problem-dependent judgement and does not belong in shared code. Omit it and Maxima asks about the sign of $a$, and the operation fails with that reason — which is correct, since the integral diverges for $a<0$.

Assume a $>$ 0

\noindent\textbf{Assumption:}\ $\displaystyle a > 0$\par

The Fourier transform of a parametrised Gaussian.

$G = \int_{-\infty}^{\infty} e^{-a t^2} e^{-i \omega t} dt$

\begin{equation}
  G = \int_{-\infty}^{\infty} e^{-a t^2} e^{-i \omega t} dt
  \label{eq:E5}
\end{equation}

Evaluate the right-hand side.

compute the right-hand side of E5

\begin{equation}
  G = {{e^ {- \left(\frac{\omega^2}{4 a}\right) } \sqrt{\pi}}\over{\sqrt{a}}}
  \label{eq:E6}
\end{equation}

Write the convolution of two Gaussians as the definition's own integral, $\int f(\tau)g(t-\tau)d\tau$.

$C = \int_{-\infty}^{\infty} e^{-\tau^2} e^{-(t-\tau)^2} d\tau$

\begin{equation}
  C = \int_{-\infty}^{\infty} e^{-\tau^2} e^{-(t-\tau)^2} d\tau
  \label{eq:E7}
\end{equation}

Evaluating gives a Gaussian of twice the variance.

compute the right-hand side of E7

\begin{equation}
  C = {{e^ {- \left(\frac{t^2}{2}\right) } \sqrt{\pi}}\over{\sqrt{2}}}
  \label{eq:E8}
\end{equation}

The Z transform of the discrete signal $(1/2)^n$. This is the kept primitive: inside is Maxima's `nusum` (closed-form infinite summation), which no general operation expresses — measured, writing `\textbackslash{}sum\_\{n=0\}\^{}\{\textbackslash{}infty\}` does not close.

z transform (1/2)\^{}n by definition

\begin{equation}
  \mathcal{Z}\{\frac{1}{2^{n}}\}(z) = \sum_{n=0}^{\infty} \frac{1}{2^{n}}\, z^{-n} = \frac{2 z}{2 z-1}-\frac{1}{2^{\infty } z^{\infty } \left(2 z-1\right)}
  \label{eq:E9}
\end{equation}

\end{document}